\documentclass{article}
\usepackage{amsmath, amssymb}
\usepackage{graphicx}
\usepackage{amsthm}
\newtheorem{conjecture}{Conjecture}
\newtheorem{remark}{Remark}
\title{Numerical Observations on Radial Scaling of $Q_{i,jk}$}
\date{}
\begin{document}
\maketitle

% Exploratory notes on the radial scaling of the sparse collision operator
% These are numerical findings from n=3 (k=0,1,2); verification at n=4 is pending.

\section*{Numerical Observations on Radial Scaling of $Q_{i,jk}$}

\subsection*{Setup and notation}

The sparse collision operator is stored as a list of $n^3$ matrices:
for each test-function index $i = (k_i, l_i, m_i)$, the matrix
$\bigl[Q_{i,st}\bigr]$ with $s = (k_s, l_s, m_s)$ and $t = (k_t, l_t, m_t)$
gives the bilinear form entries.
The basis functions are $\tilde{f}_{klm} = \mu_{k,l}\,\varphi_{klm}$, where
\[
  \mu_{k,l} = \sqrt{\frac{k!}{\Gamma(k+2l+3)}},
  \qquad
  \varphi_{klm}(r,\theta,\varphi)
    = c_{lm}\,L_k^{2l+2}(r)\,r^l\,P_l^{|m|}(\cos\theta)\,A_m(\varphi),
\]
so every entry satisfies
\[
  Q_{i,st} = \mu_{k_s,l_s}\,\mu_{k_t,l_t}\;\tilde{Q}_{i,st},
\]
where $\tilde{Q}_{i,st}$ is the integral against the un-normalised basis
$\varphi_s$ and $\varphi_t$ (the test function $\varphi_i$ carries no $\mu$
factor in the Petrov--Galerkin formulation).

\subsection*{Structural observations from the $n=3$ operator}

\paragraph{Stable support.}
For a fixed angular pair $(l_i,m_i)$, the set of nonzero positions $(s,t)$
in the operator matrix is \emph{identical} for $k_i = 0,1,2$.
In particular, the nonzero count per test-function slice does not change
with $k_i$.

\paragraph{Sign structure.}
For $l_i \geq 2$, every nonzero entry satisfies
\[
  \operatorname{sgn}\bigl(Q_{k_i,st}\bigr) = (-1)^{k_i}\,\operatorname{sgn}\bigl(Q_{0,st}\bigr),
\]
consistent with the leading sign of $L_{k_i}^{\,2l_i+2}$ through the integrand.
For $l_i = 1$ the claim does \emph{not} hold universally: some
$(k_s,l_s,k_t,l_t)$ blocks break the alternation.
This is expected because the $l_i=1$, $k_i=0$ entries
$Q_{(0,1,m),st}$ are collision-invariant projections (the $k_i=0$, $l_i=1$
test functions are the momentum-carrying basis elements), so they are
near-zero at finite quadrature order; ratios formed by dividing by these
entries are dominated by quadrature noise.
The $l_i \geq 2$ case is free of this issue.

\paragraph{$m$-independence of the scaling.}
Numerically, the ratio $Q_{k_i,st}/Q_{0,st}$ depends only on the six
integers $(k_i,l_i,k_s,l_s,k_t,l_t)$ and is \emph{independent of all
magnetic quantum numbers} $m_i,m_s,m_t$.
This is expected from the rotational structure of the spherical-harmonic
factors, which contribute only Clebsch--Gordan-type constants that cancel
in the ratio.
Verified at $k_s=3$: both $(l_s,l_t)\in\{(1,2),(2,1)\}$ give $S_3/S_0$
independent of $m_s,m_t$ to float64 precision.

\subsection*{Scaling of the stripped integral with radial degree}

Fix the test function at $(k_i = 0, l_i, m_i)$ and one field index, say
$t = (0, l_t, m_t)$, and let $k_s$ vary over $0, 1, 2$.
(The formula below is specific to $k_i = 0$; test functions with $k_i \geq 1$
carry a different $k_s$-scaling pattern that is not characterised here.)
Define the \emph{stripped} sequence
\[
  S_k \;=\; \frac{Q_{i,\,(k,\,l_s,\,m_s),\,t}}
               {\mu_{k,\,l_s}\,\mu_{0,\,l_t}},
\qquad k = 0,1,2,
\]
normalised so that $S_0 > 0$.  (An analogous sequence holds when $k_t$
varies with $k_s$ fixed; the roles of $l_s$ and $l_t$ simply swap.)

For the pairs $(l_s, l_t) \in \{(1,2),(2,1)\}$ the stripped sequence takes
\emph{exact integer} values:

\begin{center}
\begin{tabular}{cccc}
\hline
$l_s$ & $S_0, S_1, S_2$ & $S_3$ (predicted) & $S_3$ (computed) \\
\hline
$1$ & $1,\;-3,\;7$ & $-13$ & $-13$ \\
$2$ & $1,\;-5,\;16$ & $-40$ & $-40$ \\
\hline
\end{tabular}
\end{center}
The $k=3$ values were computed by calling \texttt{operator\_numba} directly
for a single tensor entry (test function $(0,1,0)$, field indices
$(3,l_s,0)$ and $(0,l_t,0)$) using the $(n_\mathrm{lag},n_\mathrm{leb})=(7,9)$
quadrature, without building the full $n=4$ tensor.
After stripping the $\mu$ factors and normalising by $S_0$, both cases
reproduce the formula to float64 precision.

\subsection*{Combinatorial pattern}

With $\alpha_s = 2l_s + 2$ the two sequences above are reproduced by
\[
  \boxed{S_k \;=\; (-1)^k\left[\binom{k+\alpha_s-2}{k}
                              + \binom{k+\alpha_s-4}{k-2}\right]}
\]
where $\binom{n}{k}=0$ whenever $k<0$ or $n<k$.
The first term is $L_k^{\alpha_s-2}(0)$ (the generalised Laguerre polynomial
at the origin); the second is a sub-leading correction that vanishes for
$k<2$.

Sequence values ($k=3$ verified; higher $k$ unverified):
\begin{align*}
  l_s = 1: &\quad S_0,S_1,S_2,\mathbf{S_3},S_4,\ldots
             = 1,-3,7,\mathbf{-13},21,\ldots \\
  l_s = 2: &\quad S_0,S_1,S_2,\mathbf{S_3},S_4,\ldots
             = 1,-5,16,\mathbf{-40},85,\ldots \\
  l_s = 3: &\quad S_0,S_1,S_2,S_3,S_4,\ldots
             = 1,-7,29,-91,238,\ldots
\end{align*}

\subsection*{Extrapolation hypothesis}

If the formula holds for $k\geq 3$, then a tensor entry at radial degree
$k_s = n-1$ (i.e.\ for the next value of $n$) can be recovered without
recomputing the quadrature:
\[
  Q_{i,\,(k_s,\,l_s,\,m_s),\,t}
  \;\approx\;
  S_{k_s}\;\mu_{k_s,\,l_s}\,\mu_{0,\,l_t}
  \;\frac{Q_{i,\,(0,\,l_s,\,m_s),\,t}}{S_0}.
\]
The same applies symmetrically for varying $k_t$.

\subsection*{Connection to the Laguerre recurrence}

The three-term recurrence
\[
  k\,L_k^\alpha(x) = (2k-1+\alpha-x)\,L_{k-1}^\alpha(x)
                   - (k-1+\alpha)\,L_{k-2}^\alpha(x)
\]
evaluated at $x=0$ (where the $x$-term drops out) gives
\[
  k\,A_k = (2k+\alpha_s-3)\,A_{k-1} - (k+\alpha_s-3)\,A_{k-2},
  \qquad \alpha_s = 2l_s+2.
\]
Because $A_k = \binom{k+\alpha_s-2}{k}$ is a binomial coefficient,
this three-term relation degenerates to the simpler \emph{first-order} recurrence
\[
  A_k = \frac{k+\alpha_s-2}{k}\,A_{k-1}, \qquad A_0 = 1,
\]
which generates the stripped sequence $S_k = (-1)^k(A_k + A_{k-2})$ cheaply
without computing factorials or Gamma functions.

\paragraph{Analytic hint from the reduction identity.}
There is a standard Laguerre identity
\[
  L_k^{\alpha-1}(x) = L_k^{\alpha}(x) - L_{k-1}^{\alpha}(x).
\]
Applying it twice at $x=0$ yields
\[
  L_k^{\alpha-2}(0) = L_k^{\alpha}(0) - 2\,L_{k-1}^{\alpha}(0) + L_{k-2}^{\alpha}(0).
\]
The basis functions carry $L_{k_s}^{2l_s+2}$, while the stripped sequence
involves $A_k = L_k^{2l_s}(0)$, a shift of $\alpha$ by $-2$.
Applying the reduction identity twice would produce such a shift with
a two-term remainder, which may explain the structure
$S_k = (-1)^k(A_k + A_{k-2})$.
The precise mechanism — whether it arises from integrations by parts on the
collision kernel or from a generating-function identity — remains to be worked out.

\subsection*{Scaling in the test-function direction ($k_i$)}

Fix the field pair $(s,t)$ completely and vary the \emph{test-function} radial
index $k_i$ for fixed $(l_i,m_i)$.
Define the magnitude ratio
\[
  R_k \;=\;
  \frac{|Q_{(k,\,l_i,\,m_i),\,s,\,t}|}{|Q_{(0,\,l_i,\,m_i),\,s,\,t}|},
  \qquad k \geq 1.
\]
The analysis below restricts to $l_i = 2$; the $l_i = 1$ case is polluted
by the near-zero collision-invariant entries discussed above.
$k_i = 0,1,2$ are read from the stored $n=3$ sparse operator;
$k_i = 3,4,5$ are computed via direct \texttt{operator\_numba} calls with
the $(n_\mathrm{lag},n_\mathrm{leb})=(7,9)$ quadrature, which reproduces
the stored values to float64 precision at $k_i=0,1,2$.

\paragraph{$m$-independence.}
$R_k$ depends only on $(k, l_i, k_s, l_s, k_t, l_t)$ and is independent of
all magnetic quantum numbers, exactly as in the $k_s$ direction.

\paragraph{No universal ratio.}
Unlike the $k_s$ formula, $R_k$ is \emph{not} the same for all $(s,t)$ blocks
with fixed $(l_s,l_t)$: it varies significantly with $k_s$ and $k_t$.
The ratio $Q_{(k,l_i,m_i),s,t}/Q_{(0,l_i,m_i),s,t}$ depends on the full
radial content $(k_s,l_s,k_t,l_t)$ of the field pair.

\paragraph{Near-constant exponent.}
For $l_i = 2$ the ratio $\log R_2 / \log R_1$ is nearly constant across all
36 $(k_s,l_s,k_t,l_t)$ blocks: mean $1.721$, std $0.030$.
This matches $\log S_2(l_i{=}2)/\log S_1(l_i{=}2) = \log 16/\log 5 = 1.7227$
to within the scatter, where $S_k$ is the same combinatorial sequence as the
$k_s$ formula (evaluated at $l_i$).  The observation is consistent with a
power-law form
\[
  R_k \;\approx\; S_k(l_i)^{\,p(s,t)},
\]
where $p(s,t)$ is block-dependent (ranging from $\approx 1.1$ to $\approx 2.0$
over the $n=3$ operator) but $k$-independent.  Because $p$ cancels in the
logarithmic ratio, the exponent $\log S_2/\log S_1$ is universal even though
$p$ is not.

\paragraph{Empirical power-law approximation in $k_i$.}
Plotting $\log R_k$ against $\log k_i$ for four representative $(l_s,l_t)$ blocks
yields \emph{approximately} straight lines across $k_i = 1,\ldots,5$.
A global log-log regression suggests the empirical approximation
\[
  R_k \;\approx\; C(s,t)\; k^{\,\alpha(s,t)},
\]
where both the prefactor $C$ and the fitted exponent $\alpha$ depend on the block.
This is not an exact model: the local log-log slope increases with $k_i$
(from $\approx 2.4$ at $k_i=1$ to $\approx 3.5$ at $k_i=5$ for the
$(k_s{=}k_t{=}0,\,l_s{=}l_t{=}2)$ block), so $\alpha$ should be understood
as an average slope over the range $k_i=1$--$5$, not an asymptotic exponent.
Whether the local slope converges to a fixed $\alpha^*$ as $k_i \to \infty$
is an open question; verifying it numerically would require computing entries
at $k_i \gg 5$, which in turn would need quadrature sizes larger than those
verified here, since $L_{k_i}^{2l_i+2}$ has $k_i$ nodes and the accuracy of
the $(n_\mathrm{lag}{=}7,\,n_\mathrm{leb}{=}9)$ quadrature for high-$k_i$
test functions has not been established.
Fitted exponents for the four blocks with $k_s = k_t = 0$ and $m_i=0$:

\begin{center}
\begin{tabular}{cc}
\hline
$(l_s,\,l_t)$ & $\alpha$ \\
\hline
$(2,2)$ & $2.88$ \\
$(0,2)$ & $2.67$ \\
$(2,0)$ & $2.55$ \\
$(1,1)$ & $2.47$ \\
\hline
\end{tabular}
\end{center}

Those four blocks all have $k_s=k_t=0$.
Extending to all nine $(k_s,k_t)\in\{0,1,2\}^2$ within the $l_s=l_t=2$ family
is consistent with the same empirical approximation and reveals how the fitted $\alpha$ depends on $(k_s,k_t)$:

\begin{center}
\begin{tabular}{c|ccc}
  $k_s \backslash k_t$ & 0 & 1 & 2 \\
  \hline
  0 & 2.88 & 3.38 & 3.86 \\
  1 & 3.38 & 3.79 & 4.18 \\
  2 & 3.86 & 4.18 & 4.51 \\
\end{tabular}
\end{center}

Two structural properties are immediate.
\emph{Symmetry:} $\alpha(k_s,k_t) = \alpha(k_t,k_s)$ exactly, consistent with
the symmetry of the bilinear operator.
\emph{Sub-additive growth:} $\alpha$ increases monotonically with both $k_s$ and
$k_t$, but the increment $\Delta\alpha$ per unit $k_s$ decreases as $k_t$ grows
(from $\approx 0.50$ at $k_t=0$ to $\approx 0.33$ at $k_t=2$), so $\alpha$ is
\emph{not} simply additive.  Specifically, $\alpha(k_s,k_t) < \alpha(k_s,0) +
\alpha(0,k_t) - \alpha(0,0)$ for all $(k_s,k_t)\neq(0,\cdot)$ and $(\cdot,0)$,
with the deficit growing to $-0.33$ at $(2,2)$.
This is in contrast to the prefactor $C = R_1$, which \emph{is} exactly additive
(see the additive structure paragraph above).

The exponents cluster in the range $\alpha \approx 2.9$--$4.5$ over the
$l_s=l_t=2$ family for $l_i=2$.
This is consistent with the large-$k$ asymptotics of $S_k(l_i)$:
the leading term $\binom{k+2l_i}{k} \sim k^{2l_i}/(2l_i)!$ gives
$S_k \sim k^4$ for $l_i=2$, so $R_k \approx S_k^p \sim k^{4p}$ asymptotically.
The fitted $\alpha$ over $k=1$--$5$ is still below the asymptotic $4p$ because
the $k^4$ regime is not yet reached for small $k$.

\paragraph{Exact additive structure for $(l_s=l_t=2)$.}
Within the angular family $l_s=l_t=2$, the $k_i=1$ ratio satisfies
\[
  R_1(k_s,k_t) \;=\; R_1(k_s,0) + R_1(0,k_t) - R_1(0,0)
\]
to float64 precision (residuals $\lesssim 10^{-11}$).
The $k_s$ and $k_t$ contributions to $R_1$ are therefore exactly additive:
knowing the edge sequence $R_1(0,k)$ suffices to reconstruct the entire
$3\times 3$ table.  Multiplied by $20$, all entries are integers:

\begin{center}
\begin{tabular}{c|ccc}
  $k_s \backslash k_t$ & 0 & 1 & 2 \\
  \hline
  0 & 220 & 292 & 355 \\
  1 & 292 & 364 & 427 \\
  2 & 355 & 427 & 490 \\
\end{tabular}
\end{center}

\paragraph{Edge-sequence recurrence.}
The increments $\Delta_k = R_1(0,k) - R_1(0,k-1)$ satisfy
\[
  \frac{\Delta_{k+1}}{\Delta_k} \;=\; \frac{2l_t+2k+1}{2l_t+2k+2},
\]
verified at $k=1$ ($\Delta_2/\Delta_1 = 7/8$ vs.\ $(2\cdot2+3)/(2\cdot2+4) = 7/8$
for $l_t=2$).  This is the ratio of a Pochhammer-type product
\[
  \Delta_k \;=\; \Delta_1\,\frac{(l_t+3/2)_{k-1}}{(l_t+2)_{k-1}},
\]
giving predicted values $\Delta_3 = 2.835$ and $R_1(0,3) = 20.585$,
verifiable with a single \texttt{operator\_numba} call at $k_i=1$, $k_t=3$.

\subsection*{Status and next steps}

\begin{itemize}
  \item $k=0,1,2$ from the stored $n=3$ sparse operator; $k=3$ computed
        directly via a single \texttt{operator\_numba} call (no full $n=4$ build).
        All four data points match the formula to float64 precision.
  \item \textbf{Comprehensive scope check:} a full scan of all 412 entries with
        $k_t=0$, $l_s\geq 1$, $k_s\in\{1,2\}$ in the $n=3$ sparse operator was
        performed.  The 104 entries with $k_i=0$ and
        $(l_s,l_t)\in\{(1,2),(2,1)\}$ all pass (every $m$ combination).
        The remaining 308 entries, which have $k_i\geq 1$ or different
        $(l_s,l_t)$ pairs, follow \emph{different} $k_s$-scaling patterns not
        covered by this formula.
  \item \textbf{Formula scope:} $k_i=0$ is required.  The $k_i$-scaling
        (test-function direction) is characterised separately in the section
        above; it does not reduce to the $k_s$ formula.
  \item \textbf{Symmetric formula verified:} varying $k_t$ with $k_s=0$ fixed
        yields the same combinatorial formula with $l_t$ replacing $l_s$
        (checked at $k_t=3$: $-13$ for $l_t=1$, $-40$ for $l_t=2$).
  \item The $l_s = 0$ case (isotropic mode) does not follow the same pattern.
  \item The analytic origin of the two-term binomial formula
        $\binom{k+\alpha-2}{k} + \binom{k+\alpha-4}{k-2}$ is unknown.
  \item Higher $k$ ($k\geq 4$) and $l_s=3$ remain unverified for the $k_s$
        formula.
  \item The $k_i$-scaling: a global log-log fit over $k_i=1$--$5$ gives an
        empirical approximation $R_k \approx C\,k^\alpha$, with fitted
        $\alpha\in[2.88,4.51]$ for $l_i=2$ over all nine $(k_s,k_t)$ blocks
        in the $l_s=l_t=2$ family.  The local log-log slope increases with
        $k_i$, so $k^\alpha$ is not an exact model.  $\alpha$ is symmetric and
        sub-additively increasing in $(k_s,k_t)$; the prefactor $C=R_1$ is
        exactly additive.  Open: extension to other $(l_s,l_t)$ families,
        analytic form of $\alpha(k_s,k_t)$, and the closed form for
        $R_1(0,0)=11$ and $\Delta_1=3.6$.
\end{itemize}

\newpage
\section*{Careful account: scaling of $Q_{(k_i,l_i,m_i),\,s,t}$ in the test-function
          radial index $k_i$}

\subsection*{What we are documenting}

We fix a pair of trial-function indices $s = (k_s,l_s,m_s)$ and
$t = (k_t,l_t,m_t)$ completely, fix the angular quantum numbers
$(l_i,m_i)$ of the test function, and ask: how does the collision operator
entry $Q_{(k_i,l_i,m_i),\,s,t}$ grow as the radial index $k_i$ increases
from $0$ to some moderate value?

This is a purely numerical question about the output of the operator
already on disk; no new tensor builds are required.  The purpose of this
section is to record the findings carefully, with all parameters stated
explicitly, so that the results can be reproduced and cited.

\subsection*{The operator visualised: sparsity structure}

Figure~\ref{fig:sparsity} shows the full sparse operator for the $n=3$
computation ($n_\mathrm{lag}=7$, $n_\mathrm{leb}=9$).  Each panel
corresponds to one test function $(k_i, l_i, m_i)$; the panels are arranged
in a $3 \times 9$ grid with rows indexed by $k_i \in \{0,1,2\}$ and columns
by the $(l_i, m_i)$ pairs in lexicographic order.
The axes of each panel are the trial-function indices $\psi_s$ (row) and
$\psi_t$ (column).

\begin{figure}[h]
  \centering
  \includegraphics[width=\textwidth]{../plot/figures/sparsity_sparse_n3_lag7_leb9.png}
  \caption{Sparse collision operator for $n=3$.  Rows: $k_i=0,1,2$.
           Columns: $(l_i,m_i)$ pairs.  Each panel entry is nonzero where
           the collision integral $Q_{(k_i,l_i,m_i),\,s,t}$ is nonzero.}
  \label{fig:sparsity}
\end{figure}

Reading down a column (fixed $(l_i,m_i)$, increasing $k_i$) reveals the
two central observations of this section.

\paragraph{Observation 1 — constant support.}
The set of nonzero positions $(s,t)$ does \emph{not} change as $k_i$
increases from $0$ to $2$.  Every panel in the same column has identical
sparsity structure.  Increasing the radial degree of the test function
neither creates new nonzero entries nor annihilates existing ones.

\paragraph{Observation 2 — alternating sign and growing magnitude.}
Although the support is fixed, the \emph{values} at the nonzero positions
are not constant.  The sign of every entry flips with each unit increment
of $k_i$ (positive at $k_i=0$, negative at $k_i=1$, positive again at
$k_i=2$), and the magnitude grows.  These two effects — sign alternation
and magnitude growth — are what the rest of this section quantifies.

\subsection*{Restricting to $\ell_i = 2$: excluding the first four columns}

The first four columns of the grid correspond to $\ell_i = 0$ (one column)
and $\ell_i = 1$ (three columns).  Even from a cursory inspection of
Figure~\ref{fig:sparsity}, these columns display different behaviour from
the rest: the $\ell_i=0$ column, for instance, does not have constant
support across the three $k_i$ rows — a structural difference that is
already visible by eye.  The $\ell_i=0,1$ test functions are associated
with the collision invariants of the relativistic Boltzmann equation and
are expected to exhibit their own peculiarities.  A careful analysis of
those columns is deferred to a later study.

Here we restrict attention entirely to $\ell_i = 2$, whose
five test functions $(k_i, 2, m_i)$ with $m_i \in \{-2,-1,0,1,2\}$ form
the red-boxed $3\times 5$ block shown in Figure~\ref{fig:sparsity_boxed}.

\begin{figure}[h]
  \centering
  \includegraphics[width=\textwidth]{../plot/figures/sparsity_sparse_n3_lag7_leb9_boxed.png}
  \caption{Same sparsity plot as Figure~\ref{fig:sparsity}, with a red box
           highlighting the $3\times 5$ block of $\ell_i=2$ test functions
           that will be studied.  The four columns to the left of the box
           ($\ell_i=0$ and $\ell_i=1$) are excluded because they correspond
           to collision-invariant projections.}
  \label{fig:sparsity_boxed}
\end{figure}

Figure~\ref{fig:sign_grid} further confirms the two observations from the
previous section, now restricted to the $\ell_i=2$ block.

\begin{figure}[h]
  \centering
  \includegraphics[width=\textwidth]{../plot/figures/ki_normalized_grid.png}
  \caption{Sign of each entry, $\operatorname{sgn} Q_{(k_i,2,m_i),s,t}
           \in \{+1,-1\}$, for the $3\times 5$ block of $\ell_i=2$ test
           functions.  Rows: $k_i\in\{0,1,2\}$; columns: $m_i\in\{-2,\ldots,2\}$.
           The identical dot pattern across all 15 panels confirms constant
           support; the complete color swap between rows $k_i=0$ and $k_i=1$
           (and back at $k_i=2$) confirms the alternating sign.}
  \label{fig:sign_grid}
\end{figure}

\subsection*{Growth of the entries with $k_i$}

Having established that the support is constant and the sign alternates,
we now look at how the \emph{magnitude} grows.
Figure~\ref{fig:growth_signed} shows the signed ratio $R_k$
(not taking absolute values) for all nonzero $(s,t)$,
so sign alternation and magnitude growth are visible simultaneously.
Figure~\ref{fig:growth_l2} shows the same data with absolute values on a
log scale, making the rate of growth easier to compare across blocks.

\begin{figure}[h]
  \centering
  \includegraphics[width=\textwidth]{../plot/figures/ki_growth_signed.png}
  \caption{Signed normalized growth
           $Q_{(k_i,2,m_i),s,t}/Q_{(0,2,m_i),s,t}$
           vs.\ $k_i \in \{0,1,2\}$, one curve per nonzero $(s,t)$,
           coloured by $(\ell_s,\ell_t)$.
           All five panels are identical, confirming $m_i$-independence.
           Every curve passes through $+1$ at $k_i=0$, dips below zero at
           $k_i=1$, and returns positive and larger at $k_i=2$.}
  \label{fig:growth_signed}
\end{figure}

A first observation is that all five panels in Figure~\ref{fig:growth_signed}
are identical to each other, confirming that the growth is independent of
the magnetic quantum number $m_i$.  Table~\ref{tab:m_indep} makes this
explicit numerically.
Fixing the block $k_s{=}k_t{=}0$, $\ell_s{=}\ell_t{=}2$ and varying $m_i$
(left subtable), or fixing the test function $\ell_i{=}2$, $m_i{=}0$ and
varying $m_s,m_t$ (right subtable), the ratios $R_1$ and $R_2$ are
identical in every row.

\begin{table}[h]
\centering
\begin{minipage}{0.50\textwidth}
\centering
{\small $m_i$-independence \quad [$k_s{=}k_t{=}0,\;\ell_s{=}\ell_t{=}2$]}
\vspace{4pt}

\begin{tabular}{ccc|rr}
\hline
$m_i$ & $m_s$ & $m_t$ & $R_1$ & $R_2$ \\
\hline
$-2$ & $-2$ & $\phantom{-}0$ & $-11$ & $56.8$ \\
$-1$ & $-2$ & $\phantom{-}1$ & $-11$ & $56.8$ \\
$\phantom{-}0$ & $-2$ & $-2$ & $-11$ & $56.8$ \\
$\phantom{-}1$ & $-2$ & $-1$ & $-11$ & $56.8$ \\
$\phantom{-}2$ & $-1$ & $-1$ & $-11$ & $56.8$ \\
\hline
\end{tabular}
\end{minipage}
\hfill
\begin{minipage}{0.45\textwidth}
\centering
{\small $m_s,m_t$-independence \quad [$\ell_i{=}2,\,m_i{=}0,\;k_s{=}k_t{=}0,\;\ell_s{=}\ell_t{=}2$]}
\vspace{4pt}

\begin{tabular}{cc|rr}
\hline
$m_s$ & $m_t$ & $R_1$ & $R_2$ \\
\hline
$-2$ & $-2$ & $-11$ & $56.8$ \\
$-1$ & $-1$ & $-11$ & $56.8$ \\
$\phantom{-}0$ & $\phantom{-}0$ & $-11$ & $56.8$ \\
$\phantom{-}1$ & $\phantom{-}1$ & $-11$ & $56.8$ \\
$\phantom{-}2$ & $\phantom{-}2$ & $-11$ & $56.8$ \\
\hline
\end{tabular}
\end{minipage}
\caption{Independence of the normalised growth ratios
         $R_k = Q_{(k_i,2,m_i),s,t}/Q_{(0,2,m_i),s,t}$
         from all magnetic quantum numbers.
         Left: $m_i$ varies, one representative $(m_s,m_t)$ per row.
         Right: $m_i=0$ fixed, $(m_s,m_t)$ varies.
         All entries give $R_1=-11$ and $R_2=56.8$ exactly.}
\label{tab:m_indep}
\end{table}

\begin{conjecture}[Independence from magnetic quantum numbers]
For fixed radial and angular indices $(k_i, \ell_i)$, $(k_s, \ell_s)$,
$(k_t, \ell_t)$, the normalised growth ratio
\[
  R_k \;=\; \frac{Q_{(k,\,\ell_i,\,m_i),\,s,\,t}}{Q_{(0,\,\ell_i,\,m_i),\,s,\,t}}
\]
is independent of all magnetic quantum numbers $m_i$, $m_s$, $m_t$,
whenever the entry is nonzero.
\end{conjecture}

\begin{remark}
The two parts of the conjecture have different levels of directness.

\emph{$m_s, m_t$-independence (Table~\ref{tab:m_indep}, right).}
This is the cleaner case.  For a fixed test function $(k_i, \ell_i, m_i)$,
the operator slice is a single matrix.  Within that matrix, all entries
sharing the same $(k_s, \ell_s, k_t, \ell_t)$ but differing in $(m_s, m_t)$
are directly comparable nonzero entries of the \emph{same} slice.
Table~\ref{tab:m_indep} (right) shows five such entries all giving
$R_1 = -11$, $R_2 = 56.8$.

\emph{$m_i$-independence (Table~\ref{tab:m_indep}, left).}
This is more subtle.  The support of the operator slice — which $(s,t)$
entries are nonzero — itself shifts with $m_i$ through the angular selection
rules, so one cannot fix a single $(m_s, m_t)$ pair and compare the same
entry across different $m_i$ values.  The claim is that at whichever
positions a nonzero entry \emph{does} appear, the ratio takes the same value
regardless of the particular $(m_i, m_s, m_t)$ combination.
Table~\ref{tab:m_indep} (left) shows five entries from five different slices,
each at a different matrix position, all giving the same $R_1$ and $R_2$.
\end{remark}

\begin{remark}[Why $m$-independence is expected]
The basis functions factor as
\[
  \varphi_{k\ell m}(r,\theta,\varphi)
  \;=\;
  \underbrace{c_{\ell m}\,P_\ell^{|m|}(\cos\theta)\,A_m(\varphi)}_{\text{angular: depends on }\ell,m}
  \;\times\;
  \underbrace{L_k^{2\ell+2}(r)\,r^\ell}_{\text{radial: depends on }k,\ell\text{ only}}.
\]
Increasing $k$ changes \emph{only} the radial Laguerre polynomial
$L_k^{2\ell+2}$; the spherical-harmonic factor, which carries all the $m$
dependence, is unaffected.
In the ratio $R_k$, the angular factors of the test function appear in both
numerator and denominator and cancel exactly, leaving a quantity that
depends only on the radial indices $(k, \ell_i, k_s, \ell_s, k_t, \ell_t)$.
\end{remark}

\begin{figure}[h]
  \centering
  \includegraphics[width=\textwidth]{../plot/figures/ki_growth_l2.png}
  \caption{Absolute normalized growth
           $|Q_{(k_i,2,m_i),s,t}|/|Q_{(0,2,m_i),s,t}|$
           on a log scale.  The spread between curves reflects the
           dependence of the growth rate on $(\ell_s,\ell_t)$.}
  \label{fig:growth_l2}
\end{figure}

\subsection*{Does $A_k$ satisfy a closed recurrence in $k_i$?}

Since the collision operator is linear in the test function and the generalised
Laguerre polynomials satisfy the three-term recurrence
\[
  (k+1)\,L_{k+1}^{\alpha}(r)
  = (2k+\alpha+1-r)\,L_k^{\alpha}(r)
  - (k+\alpha)\,L_{k-1}^{\alpha}(r),
\]
a natural question is whether the sequence
$A_k = Q_{(k,\,\ell_i,\,m_i),\,s,\,t}\,/\,Q_{(0,\,\ell_i,\,m_i),\,s,\,t}$
inherits a similar recurrence.
We tested two forms using the representative entry
$(\ell_i{=}2,\,m_i{=}0)$, $s{=}t{=}(0,2,0)$.

\paragraph{Form~(I): Laguerre-type, 2 parameters.}
\[
  (k+1)\,A_{k+1} = (2k+\beta)\,A_k - (k+\gamma)\,A_{k-1}.
\]
Fitting $\beta$ and $\gamma$ from $k=1,2$ (two equations, two unknowns) and
predicting $k=3,4,5$ gives relative errors of $3.7\%$, $11.7\%$, and $22.5\%$
respectively.  The 2-parameter Laguerre form does \emph{not} hold.
The reason is clear in hindsight: applying the Laguerre recurrence to the
test function introduces a term $\int r'\,L_k^{\alpha}(r')\cdot[\text{integrand}]$,
where $r'$ is the post-collision momentum magnitude varying over the quadrature.
Because $r'$ is not a single effective value, the $r'$-correction does not reduce
to a multiple of $A_k$ or $A_{k-1}$, and the recurrence does not close at 2 parameters.

\paragraph{Form~(II): general linear, 4 parameters.}
\[
  (k+1)\,A_{k+1} = (ak+b)\,A_k - (ck+d)\,A_{k-1}.
\]
Fitting $a,b,c,d$ exactly from the four equations at $k=1,2,3,4$ and predicting
$A_5$ gives a relative error of $\lesssim 10^{-12}$ (float64 rounding only),
which initially suggested the sequence might be P-recursive.
However, computing $A_6$ refutes this: the prediction overshoots by $2.7\%$.
The value $A_6 \approx 2098.4$ was computed independently with two quadrature
sizes, $(n_\mathrm{lag},n_\mathrm{leb})=(7,9)$ and $(11,13)$, and the two
agree to $2\times 10^{-12}$ relative error, confirming the discrepancy is not
a quadrature artefact.
(The $(11,13)$ quadrature file on disk had been truncated; it was rebuilt
using a numpy-vectorised version of the quadrature loop and verified to
reproduce the $(7,9)$ operator entries to float64 precision before use.)

\paragraph{Conclusion: the sequence is not order-2 P-recursive.}
The 4-parameter degree-1 recurrence does not hold beyond $k=5$.
The apparent success at $k=5$ was based on a single free prediction and
turned out to be coincidental.
The $k_i$ direction is fundamentally harder than the $k_s$ direction:
in the $k_s$ formula the Laguerre polynomial was evaluated at $r=0$, collapsing
to pure combinatorics; in the $k_i$ direction the test-function Laguerre
polynomial $L_{k_i}^{2\ell_i+2}(r')$ is evaluated at the post-collision
momentum $r'$, which varies continuously over the quadrature and admits no
such simplification.

\paragraph{What the data does show.}
The seven computed values
\[
  (A_0,\ldots,A_6)
  = \tfrac{1}{5}(5,\,-55,\,284,\,-980,\,2602,\,-5678,\,10492)
\]
are all rational with denominator~$5$ (or integer), and the unsigned
magnitudes grow monotonically.
The local log-log slopes
\[
  \alpha_k = \frac{\log|A_k/A_{k-1}|}{\log(k/(k-1))}
\]
are $2.37,\,3.06,\,3.39,\,3.50,\,3.38$ for $k=2,\ldots,6$,
suggesting a possible convergence toward $\alpha^* \approx 3.4$--$3.5$
rather than continuing to grow.
Whether the sequence has a closed-form pattern, and what asymptotic
power law (if any) it approaches, remain open questions.

\end{document}
